% ============================================================
% 辐射章重构 part3(对应原 script.tex 第 1070--1729 行)
% 三个 \subsubsection 属于 \subsection{点粒子的辐射反作用}
% 风格:公式全部无编号、带量纲量、不引入 K/U 等未声明简写
% ============================================================

\subsubsection{世界管几何与管壁四动量通量}

上一小节只证明了正则场 $F_{\mathrm{reg}}$ 在世界线上有有限极限;这还不足以确定运动方程,因为质量重整化来自奇异场携带的四动量,而不是世界线上的场值。为此,在世界线周围取固有半径(管壁到世界线的固有距离)为 $\ell$ 的世界管,计算奇异场穿过管壁的四动量通量。管壁参数化为

\begin{equation*}
 y^\mu(\tau,\theta,\phi)
 =z^\mu(\tau)+\ell n^\mu(\tau,\theta,\phi),
 \qquad
 n^2=-1,
 \qquad
 n\cdot u=0.
\end{equation*}

在每个 $\tau$ 的瞬时静止系中,

\begin{equation*}
 u^\mu=(c,\vb{0}),
 \qquad
 n^\mu=(0,\sin\theta\cos\phi,\sin\theta\sin\phi,\cos\theta).
\end{equation*}

三个切向量为

\begin{equation*}
 {\partial y^\mu\over\partial\tau}=u^\mu+\ell\dot n^\mu,
 \qquad
 {\partial y^\mu\over\partial\theta}=\ell\partial_\theta n^\mu,
 \qquad
 {\partial y^\mu\over\partial\phi}=\ell\partial_\phi n^\mu.
\end{equation*}

管壁有向面积元定义为(设 $\epsilon_{0123}=+1$)

\begin{equation*}
 \d\Sigma_\nu
 =\epsilon_{\nu\alpha\beta\xi}
 {\partial y^\alpha\over\partial\tau}
 {\partial y^\beta\over\partial\theta}
 {\partial y^\xi\over\partial\phi}
 \d\tau\d\theta\d\phi.
\end{equation*}

由 $n\cdot u=0$ 对固有时微分得

\begin{equation*}
 \dot n\cdot u=-n\cdot a.
\end{equation*}

在瞬时静止系中,切向量矩阵的 $3\times3$ 非零子式为 $\ell^2\sin\theta\,(c+\ell\dot n^0)$:空间部分 $\partial_\theta n^i$、$\partial_\phi n^i$ 张成切平面,其混合积 $|\partial_\theta n\times\partial_\phi n|=\sin\theta$;由上式得 $\dot n^0=-n\cdot a/c$,故行列式为

\begin{equation*}
 \ell^2\sin\theta\left(c+\ell\dot n^0\right)
 =c\ell^2\sin\theta\left(1-{\ell n\cdot a\over c^2}\right).
\end{equation*}

外向空间法矢为 $n^\mu$,但降指标后 $n_i=-n^i$,故

\begin{equation*}
 \d\Sigma_\nu
 =-c\ell^2\left(1-{\ell n\cdot a\over c^2}\right)n_\nu\d\tau\d\Omega
 =-c\ell\left(1-{\xi\cdot a\over c^2}\right)\xi_\nu\d\tau\d\Omega,
\end{equation*}

最后一步用了 $\xi^\mu=\ell n^\mu$。这里出现的几何因子正是近线展开中的 $\Lambda=1-\xi\cdot a/c^2$。

在号差 $+2$ 下,SI 制电磁能动量张量为

\begin{equation*}
 T_{\mathrm{em}}^{\mu\nu}
 ={1\over\mu_0}
 \left(
 -F^{\mu\alpha}F^\nu{}_{\alpha}
 +{1\over4}g^{\mu\nu}F^{\alpha\beta}F_{\alpha\beta}
 \right).
\end{equation*}

管壁上的四动量通量为

\begin{equation*}
 \d P_{\mathrm{wall}}^\mu
 ={1\over c}T_{\mathrm{em}}^{\mu\nu}\d\Sigma_\nu.
\end{equation*}

系数 $1/c$ 来自 SI 制中 $x^0=ct$:四动量的空间分量为 $\int T^{i0}\d^3x/c$,时间分量为 $\int T^{00}\d^3x/c=E/c$。由于面积元为 $O(\ell^2)$,若希望保留积分后发散的 $O(\ell^{-1})$ 项和有限的 $O(1)$ 项,能动量张量中必须保留奇异场平方到 $O(\ell^{-3})$,以及奇异场与正则场的交叉项到 $O(\ell^{-2})$。

为缩短中间代数,定义三个组合

\begin{equation*}
 \mathcal G^{\mu\nu}={\xi^{[\mu}u^{\nu]}\over c},
 \qquad
 \mathcal H^{\mu\nu}={\xi^{[\mu}\dot a^{\nu]}\over c^3},
 \qquad
 \mathcal J^{\mu\nu}={u^{[\mu}a^{\nu]}\over c^3},
\end{equation*}

它们分别承载奇异场展开中的 $\xi u$、$\xi\dot a$ 与 $u a$ 三类反对称结构。上一小节的奇异场展开式使实际场在管壁上写成

\begin{equation*}
 F_{\mathrm{act}}^{\mu\nu}
 =F_{\mathrm{reg}}^{\mu\nu}
 +{q\over4\pi\varepsilon_0c}\Lambda^{-1/2}
 \left[
 A\mathcal G^{\mu\nu}
 +B\mathcal J^{\mu\nu}
 +C\mathcal H^{\mu\nu}
 \right]+O(\ell),
\end{equation*}

其中

\begin{equation*}
 A=\ell^{-3}-{a^2\over8c^4}\ell^{-1},
 \qquad
 B={1+\xi\cdot a/c^2\over2}\ell^{-1},
 \qquad
 C=-{1\over2}\ell^{-1}.
\end{equation*}

计算张量平方所需的乘积先保持未收缩形式。由反对称化定义逐项相乘,并用 $\xi\cdot u=0$、$u^2=c^2$、$u\cdot a=0$、$u\cdot\dot a=-a^2$、$\xi^2=-\ell^2$,得到

\begin{align*}
 \mathcal G^\mu{}_\alpha\mathcal G^{\alpha\nu}
 &=-\xi^\mu\xi^\nu+\ell^2{u^\mu u^\nu\over c^2},\\
 \mathcal G^\mu{}_\alpha\mathcal J^{\alpha\nu}
 &={\xi^\mu a^\nu\over c^2}+{(\xi\cdot a)u^\mu u^\nu\over c^4},\\
 \mathcal J^\mu{}_\alpha\mathcal G^{\alpha\nu}
 &={a^\mu\xi^\nu\over c^2}+{(\xi\cdot a)u^\mu u^\nu\over c^4},\\
 \mathcal G^\mu{}_\alpha\mathcal H^{\alpha\nu}
 &={a^2\xi^\mu\xi^\nu\over c^4}
 +{\ell^2u^\mu\dot a^\nu\over c^4}
 +{(\xi\cdot\dot a)u^\mu\xi^\nu\over c^4},\\
 \mathcal H^\mu{}_\alpha\mathcal G^{\alpha\nu}
 &={a^2\xi^\mu\xi^\nu\over c^4}
 +{\ell^2\dot a^\mu u^\nu\over c^4}
 +{(\xi\cdot\dot a)\xi^\mu u^\nu\over c^4},\\
 \mathcal J^\mu{}_\alpha\mathcal J^{\alpha\nu}
 &=-{a^2u^\mu u^\nu\over c^6}-{a^\mu a^\nu\over c^4}.
\end{align*}

例如第一式逐项展开为

\begin{align*}
 {1\over c^2}(\xi^\mu u_\alpha-u^\mu\xi_\alpha)
 (\xi^\alpha u^\nu-u^\alpha\xi^\nu)
 &={(\xi\cdot u)\xi^\mu u^\nu\over c^2}
 -{u^2\xi^\mu\xi^\nu\over c^2}
 -{\xi^2u^\mu u^\nu\over c^2}
 +{(\xi\cdot u)u^\mu\xi^\nu\over c^2}\\
 &=-\xi^\mu\xi^\nu+{\ell^2u^\mu u^\nu\over c^2};
\end{align*}

其余各式完全同理,例如 $\mathcal G^\mu{}_\alpha\mathcal J^{\alpha\nu}$ 中 $(\xi\cdot u)u^\mu a^\nu/c^4$、$\xi^\mu(u\cdot a)u^\nu/c^4$、$u^\mu(\xi\cdot u)a^\nu/c^4$ 三项分别因 $\xi\cdot u=0$ 与 $u\cdot a=0$ 消失,只留下 $\xi^\mu a^\nu/c^2$ 与 $(\xi\cdot a)u^\mu u^\nu/c^4$。

先求能动量张量第一项中出现的

\begin{equation*}
 \mathcal X^\mu
 \equiv F_{\mathrm{act}}^\mu{}_\alpha
 F_{\mathrm{act}}^{\alpha\nu}\xi_\nu.
\end{equation*}

把管壁上的场展开为奇异场平方与奇异--正则交叉项:

\begin{align*}
 \mathcal X^\mu
 ={}&{q^2\over16\pi^2\varepsilon_0^2c^2\Lambda}
 \bigg[
 A^2\mathcal G^\mu{}_\alpha\mathcal G^{\alpha\nu}
 +AB\left(
 \mathcal G^\mu{}_\alpha\mathcal J^{\alpha\nu}
 +\mathcal J^\mu{}_\alpha\mathcal G^{\alpha\nu}
 \right)\\
 &+AC\left(
 \mathcal G^\mu{}_\alpha\mathcal H^{\alpha\nu}
 +\mathcal H^\mu{}_\alpha\mathcal G^{\alpha\nu}
 \right)
 +B^2\mathcal J^\mu{}_\alpha\mathcal J^{\alpha\nu}
 \bigg]\xi_\nu\\
 &+{q\over4\pi\varepsilon_0c}\ell^{-3}
 \left[
 \mathcal G^\mu{}_\alpha F_{\mathrm{reg}}^{\alpha\nu}
 +F_{\mathrm{reg}}^\mu{}_\alpha\mathcal G^{\alpha\nu}
 \right]\xi_\nu
 +O(1).
\end{align*}

这里 $BH$、$CH$ 和 $H^2$ 至多在面积积分后产生 $O(\ell)$,所以没有写出;交叉项中也只需保留最强的 $A\mathcal G\sim\ell^{-2}$,其 $\Lambda^{-1/2}$ 因子已并入 $\ell^{-3}$ 的 $O(\ell^2)$ 修正中。

把乘积公式与 $\xi_\nu$ 收缩。例如

\begin{align*}
 (\mathcal G\mathcal G)^{\mu\nu}\xi_\nu
 &=\left(-\xi^\mu\xi^\nu+{\ell^2u^\mu u^\nu\over c^2}\right)\xi_\nu
 =-\xi^\mu(\xi\cdot\xi)+{\ell^2u^\mu(\xi\cdot u)\over c^2}
 =\ell^2\xi^\mu,\\
 (\mathcal G\mathcal J)^{\mu\nu}\xi_\nu
 &=\xi^\mu{(a\cdot\xi)\over c^2}+{(\xi\cdot a)u^\mu(u\cdot\xi)\over c^4}
 ={(\xi\cdot a)\xi^\mu\over c^2},
\end{align*}

其余各式用同一套恒等式,完整收缩表为

\begin{align*}
 (\mathcal G\mathcal G)^{\mu\nu}\xi_\nu
 &=\ell^2\xi^\mu,\\
 (\mathcal G\mathcal J)^{\mu\nu}\xi_\nu
 &={(\xi\cdot a)\xi^\mu\over c^2},\\
 (\mathcal J\mathcal G)^{\mu\nu}\xi_\nu
 &=-{\ell^2a^\mu\over c^2},\\
 (\mathcal G\mathcal H)^{\mu\nu}\xi_\nu
 &=-{\ell^2a^2\xi^\mu\over c^4},\\
 (\mathcal H\mathcal G)^{\mu\nu}\xi_\nu
 &=-{\ell^2a^2\xi^\mu\over c^4},\\
 (\mathcal J\mathcal J)^{\mu\nu}\xi_\nu
 &=-{(\xi\cdot a)a^\mu\over c^4}.
\end{align*}

正则场交叉项则为

\begin{align*}
 &\left[
 \mathcal G^\mu{}_\alpha F_{\mathrm{reg}}^{\alpha\nu}
 +F_{\mathrm{reg}}^\mu{}_\alpha\mathcal G^{\alpha\nu}
 \right]\xi_\nu\\
 &\qquad
 ={\ell^2\over c}F_{\mathrm{reg}}^\mu{}_\nu u^\nu
 +{\xi^\mu\over c}F_{\mathrm{reg}\,\alpha\beta}u^\alpha\xi^\beta,
\end{align*}

其中第一项来自 $\mathcal G^{\alpha\nu}\xi_\nu=\ell^2u^\alpha/c$,第二项来自 $\mathcal G^\mu{}_\alpha F_{\mathrm{reg}}^{\alpha\nu}\xi_\nu$,而 $F_{\mathrm{reg}}^{\alpha\nu}\xi_\nu\xi_\alpha=0$(反对称张量收缩于对称张量)。

将 $A,B,C$ 代入,并用到

\begin{align*}
 A^2&=\ell^{-6}-{a^2\over4c^4}\ell^{-4}+O(\ell^{-2}),\\
 AB&={1+\xi\cdot a/c^2\over2}\ell^{-4}+O(\ell^{-2}),\\
 AC&=-{1\over2}\ell^{-4}+O(\ell^{-2}),\\
 B^2&={1\over4}\ell^{-2}+O(\ell^{-1}),
\end{align*}

$\mathcal X^\mu$ 逐项化成

\begin{align*}
 \mathcal X^\mu
 ={q^2\over16\pi^2\varepsilon_0^2c^2\Lambda}
 \bigg\{{}
 &\left(\ell^{-6}-{a^2\over4c^4}\ell^{-4}\right)
 \ell^2\xi^\mu\\
 &+{1+\xi\cdot a/c^2\over2}\ell^{-4}
 \left[{(\xi\cdot a)\xi^\mu\over c^2}
 -{\ell^2a^\mu\over c^2}\right]\\
 &+{a^2\ell^{-2}\xi^\mu\over c^4}
 -{(\xi\cdot a)a^\mu\ell^{-2}\over4c^4}
 \bigg\}\\
 &+{q\over4\pi\varepsilon_0c}\ell^{-3}
 \left[
 {\ell^2\over c}F_{\mathrm{reg}}^\mu{}_\nu u^\nu
 +{\xi^\mu\over c}F_{\mathrm{reg}\,\alpha\beta}u^\alpha\xi^\beta
 \right]+O(1).
\end{align*}

这就是 $F^\mu{}_\alpha F^{\alpha\nu}\xi_\nu$ 的完整所需展开。

再求能动量张量的迹项。所需的反对称张量全收缩为

\begin{align*}
 \mathcal G_{\alpha\beta}\mathcal G^{\alpha\beta}
 &=-2\ell^2,\\
 \mathcal G_{\alpha\beta}\mathcal J^{\alpha\beta}
 &=-{2\xi\cdot a\over c^2},\\
 \mathcal G_{\alpha\beta}\mathcal H^{\alpha\beta}
 &={2\ell^2a^2\over c^4},\\
 \mathcal J_{\alpha\beta}\mathcal J^{\alpha\beta}
 &={2a^2\over c^4}.
\end{align*}

例如

\begin{align*}
 \mathcal G_{\alpha\beta}\mathcal J^{\alpha\beta}
 &={1\over c^4}(\xi_\alpha u_\beta-\xi_\beta u_\alpha)
 (u^\alpha a^\beta-u^\beta a^\alpha)\\
 &={1\over c^4}\left[-(\xi\cdot a)u^2-(u\cdot u)(\xi\cdot a)\right]
 =-{2\xi\cdot a\over c^2},
\end{align*}

其中利用了 $u\cdot a=0$、$u^2=c^2$。因此

\begin{align*}
 F_{\mathrm{act}\,\alpha\beta}F_{\mathrm{act}}^{\alpha\beta}
 ={}&{q^2\over16\pi^2\varepsilon_0^2c^2\Lambda}
 \bigg[
 -2A^2\ell^2
 -4AB{\xi\cdot a\over c^2}
 +4AC{\ell^2a^2\over c^4}
 +2B^2{a^2\over c^4}
 \bigg]\\
 &+{q\over\pi\varepsilon_0c}\ell^{-3}
 {F_{\mathrm{reg}\,\alpha\beta}\xi^\alpha u^\beta\over c}
 +O(1),
\end{align*}

其中 $-4AB(\xi\cdot a)$ 来自交叉项 $2AB\bigl[\mathcal G_{\alpha\beta}\mathcal J^{\alpha\beta}+\mathcal J_{\alpha\beta}\mathcal G^{\alpha\beta}\bigr]=2AB\cdot2(-2\xi\cdot a/c^2)$ 型收缩的两重组合。把 $A,B,C$ 代入并保留所需阶数,得到

\begin{align*}
 F_{\mathrm{act}\,\alpha\beta}F_{\mathrm{act}}^{\alpha\beta}
 ={}&{q^2\over16\pi^2\varepsilon_0^2c^2\Lambda}
 \left[
 -2\left(1+{\xi\cdot a\over c^2}
 +\left({\xi\cdot a\over c^2}\right)^2\right)\ell^{-4}
 -{a^2\ell^{-2}\over c^4}
 \right]\\
 &+{q\over\pi\varepsilon_0c}\ell^{-3}
 {F_{\mathrm{reg}\,\alpha\beta}\xi^\alpha u^\beta\over c}
 +O(1).
\end{align*}

于是迹项为

\begin{align*}
 {1\over4}\xi^\mu
 F_{\mathrm{act}\,\alpha\beta}F_{\mathrm{act}}^{\alpha\beta}
 ={}&{q^2\over16\pi^2\varepsilon_0^2c^2\Lambda}
 \left[
 -{1\over2}
 \left(1+{\xi\cdot a\over c^2}
 +\left({\xi\cdot a\over c^2}\right)^2\right)\ell^{-4}
 -{a^2\ell^{-2}\over4c^4}
 \right]\xi^\mu\\
 &+{q\over4\pi\varepsilon_0c}\ell^{-3}
 {\xi^\mu F_{\mathrm{reg}\,\alpha\beta}\xi^\alpha u^\beta\over c}
 +O(1).
\end{align*}

现在把 $\mathcal X^\mu$ 的展开与迹项相加。由于 $F_{\mathrm{reg}\,\alpha\beta}$ 反对称,

\begin{equation*}
 F_{\mathrm{reg}\,\alpha\beta}u^\alpha\xi^\beta
 =-F_{\mathrm{reg}\,\alpha\beta}\xi^\alpha u^\beta,
\end{equation*}

两个正则场交叉项中与 $\xi^\mu$ 成正比的部分相消,只留下 $F_{\mathrm{reg}}^\mu{}_\nu u^\nu$;$\xi^\mu$ 的系数按 $\ell$ 幂次收集:$\ell^{-4}$ 部分为 $1+\frac12(\xi\cdot a/c^2)+\frac12(\xi\cdot a/c^2)^2-\frac12\bigl(1+\xi\cdot a/c^2+(\xi\cdot a/c^2)^2\bigr)=\frac12$,平方项相消,$\ell^{-2}$ 部分为 $-a^2/(4c^4)+a^2/c^4-a^2/(4c^4)=a^2/(2c^4)$;$a^\mu$ 的系数为 $-(1+\xi\cdot a/c^2)/(2c^2)-(\xi\cdot a/c^2)/(4c^2)=-(1+3\xi\cdot a/2c^2)/(2c^2)$。同一组合化成简洁形式

\begin{align*}
 \mathcal Q^\mu
 &\equiv
 F_{\mathrm{act}}^\mu{}_\alpha
 F_{\mathrm{act}}^{\alpha\nu}\xi_\nu
 +{1\over4}\xi^\mu
 F_{\mathrm{act}\,\alpha\beta}F_{\mathrm{act}}^{\alpha\beta}\\
 &={q^2\over16\pi^2\varepsilon_0^2c^2\Lambda}
 \left[
 \left({1\over2}\ell^{-4}
 +{a^2\ell^{-2}\over2c^4}\right)\xi^\mu
 -{\ell^{-2}\over2c^2}
 \left(1+{3\xi\cdot a\over2c^2}\right)a^\mu
 \right]\\
 &\quad
 +{q\over4\pi\varepsilon_0c^2}\ell^{-1}
 F_{\mathrm{reg}}^\mu{}_\nu u^\nu
 +O(1).
\end{align*}

这一步把看似众多的张量平方化成了三个结构:沿径向的奇函数项、沿加速度的发散惯性项,以及有限的正则 Lorentz 力。

把能动量张量收缩到面积元上:$\d P^\mu=(1/c)T_{\mathrm{em}}^{\mu\nu}\d\Sigma_\nu$ 中,奇异场平方部分恰好按 $\mathcal Q^\mu$ 组装,其中面积元因子 $(1-\xi\cdot a/c^2)$ 与 $\mathcal Q^\mu$ 携带的 $\Lambda^{-1}$ 精确合并为 $1$(因为 $\Lambda=1-\xi\cdot a/c^2$);正则交叉部分合并面积元因子后给出与 $qF_{\mathrm{reg}}^{\mu\nu}u_\nu$ 成正比的项。使用 $\mu_0\varepsilon_0c^2=1$,单位固有时和单位立体角的管壁四动量通量为

\begin{align*}
 {\d^2P_{\mathrm{wall}}^\mu\over\d\tau\d\Omega}
 =-{1\over4\pi}\bigg\{ {q^2\over4\pi\varepsilon_0}
 \bigg[{}
 &\left({1\over2\ell^3}
 +{a^2\over2c^4\ell}\right)\xi^\mu
 -{a^\mu\over2c^2\ell}
 -{3(\xi\cdot a)a^\mu\over4c^4\ell}
 \bigg]\\
 &+\left(1+{\xi\cdot a\over c^2}\right)
 qF_{\mathrm{reg}}^{\mu\nu}u_\nu
 \bigg\}+O(\ell).
\end{align*}

这是同一个通量,但每项的角奇偶性已经显式:方括号内第一、二项为径向奇函数,第三、四项沿加速度方向。

瞬时静止系中的球面对称给出

\begin{equation*}
 \int n^\mu\d\Omega=0,
 \qquad
 \int n^\mu n^\nu\d\Omega
 =-{4\pi\over3}
 \left(g^{\mu\nu}-{u^\mu u^\nu\over c^2}\right).
\end{equation*}

因为 $\xi^\mu=\ell n^\mu$,通量中所有含单个 $\xi^\mu$ 的项积分为零;含 $(\xi\cdot a)a^\mu$ 的项也因 $\int\xi\cdot a\,\d\Omega=0$ 而消失。正则场在 $\ell\to0$ 时与角方向无关。角积分后只剩两项:

\begin{align*}
 \int\left[-{1\over4\pi}\cdot{q^2\over4\pi\varepsilon_0}
 \left(-{a^\mu\over2c^2\ell}\right)\right]\d\Omega
 &={q^2a^\mu\over8\pi\varepsilon_0c^2\ell},\\
 \int\left[-{1\over4\pi}
 \left(1+{\xi\cdot a\over c^2}\right)
 qF_{\mathrm{reg}}^{\mu\nu}u_\nu\right]\d\Omega
 &=-qF_{\mathrm{reg}}^{\mu\nu}u_\nu,
\end{align*}

故同一个管壁通量化成

\begin{equation*}
 {\d P_{\mathrm{wall}}^\mu\over\d\tau}
 ={q^2\over8\pi\varepsilon_0c^2\ell}a^\mu
 -qF_{\mathrm{reg}}^{\mu\nu}u_\nu
 +O(\ell).
\end{equation*}

这就是整个近线展开和世界管积分后的简洁结果:第一项与加速度平行,表现为随管半径发散的局部惯性,其系数 $q^2/(8\pi\varepsilon_0c^2\ell)$ 正是电磁质量;第二项是在世界线上有限的正则 Lorentz 四力。下一步把世界管看成封闭四维区域的边界,把管壁通量与端帽上的动量变化联系起来。

\subsubsection{端帽动量、质量重整化与有限自力}

世界管本身并不是封闭曲面:它有两个端帽。取由管壁以及 $\tau_1$、$\tau_2$ 两个端帽围成的四维区域 $\mathcal V$。真空区域中场能动量的 Gauss 定理为

\begin{equation*}
 \int_{\mathcal V}\partial_\nu T_{\mathrm{em}}^{\mu\nu}\d^4x
 =\int_{\partial\mathcal V}T_{\mathrm{em}}^{\mu\nu}\d\Sigma_\nu
 =\Delta P_{\mathrm{cap}}^\mu+P_{\mathrm{wall}}^\mu.
\end{equation*}

如果只讨论管外的真空区域,$\partial_\nu T_{\mathrm{em}}^{\mu\nu}=0$,于是管壁流出的四动量必须由两个端帽所截获的管内四动量变化补偿。不过,点粒子世界线本身已从积分区域中挖去;Maxwell 方程并不单独规定被挖去部分含有哪些非电磁自由度。因此还必须加入一个点粒子局域性假设:管内未解析自由度可由只依赖当前世界线状态的端帽四动量 $p_0^\mu$ 描述。

若粒子没有额外内部方向,并且在最低导数阶只允许使用 $u^\mu$,则

\begin{equation*}
 p_0^\mu(\ell)=m_0(\ell)u^\mu.
\end{equation*}

这正是 Dirac 把管壁通量视为端点量之差、并选择 $B^\mu\propto u^\mu$ 的物理内容。若允许 $a^\mu$ 或更高导数进入 $p_0^\mu$,便等价于给点粒子增加新的内部结构常数;那将是另一个有效理论,而不是最小的点电子模型。

在无外场时,端帽内粒子动量与管壁流量满足

\begin{equation*}
 {\d p_0^\mu\over\d\tau}
 +{\d P_{\mathrm{wall}}^\mu\over\d\tau}=0.
\end{equation*}

有外部无源场时,正则场已经包含 $F_{\mathrm{in}}^{\mu\nu}$,故把管壁通量公式代入,并使用 ${\d p_0^\mu/\d\tau}=m_0(\ell)a^\mu$,得到预重整化方程

\begin{equation*}
 qF_{\mathrm{reg}}^{\mu\nu}u_\nu
 =\left[
 m_0(\ell)
 +{q^2\over8\pi\varepsilon_0c^2\ell}
 \right]a^\mu+O(\ell).
\end{equation*}

括号内的两个量分别发散,但它们只以同一个加速度系数出现。实验所测得的惯性质量因而定义为这一组合的有限极限:

\begin{equation*}
 m
 =\lim_{\ell\to0}
 \left[
 m_0(\ell)
 +{q^2\over8\pi\varepsilon_0c^2\ell}
 \right].
\end{equation*}

这不是把无限量"丢掉",而是说明裸质量 $m_0(\ell)$ 本身不是可测量参数;只有这一组合进入低能运动方程。

此处的发散系数

\begin{equation*}
 {q^2\over8\pi\varepsilon_0c^2\ell}
 ={1\over c^2}{q^2\over8\pi\varepsilon_0\ell}
 ={U_{\mathrm C}(\ell)\over c^2}
\end{equation*}

等于半径 $\ell$ 的 Coulomb 外场能 $U_{\mathrm C}(\ell)=q^2/(8\pi\varepsilon_0\ell)$ 除以 $c^2$,即球对称近场的 $P_{\mathrm{bound}}^0/c$ 型质量。它来自与世界线正交的协变端帽极限;前面实验室同时面上扩展球电荷的 $4U_{\mathrm{es}}/(3c^2)$ 来自非封闭电磁子系统的动量积分。两者对应不同的积分超曲面和不同的系统划分,不能仅凭"电磁质量"这一名称互相替换。

取预重整化方程的有限极限,得到最紧凑的重整化方程

\begin{equation*}
 ma^\mu=qF_{\mathrm{reg}}^{\mu\nu}u_\nu.
\end{equation*}

这里的 $F_{\mathrm{reg}}$ 仍是 Green 函数一节定义的同一个物理量。把它重新展开为

\begin{equation*}
 F_{\mathrm{reg}}^{\mu\nu}
 =F_{\mathrm{in}}^{\mu\nu}+F_{\mathrm R}^{\mu\nu},
\end{equation*}

并代入前一节的正则场在世界线上的有限极限式 $F_{\mathrm R}^{\mu\nu}(z)=(q/6\pi\varepsilon_0c^5)(\dot a^\mu u^\nu-\dot a^\nu u^\mu)$,有限自场四力为

\begin{align*}
 f_{\mathrm{self}}^\mu
 &\equiv qF_{\mathrm R}^{\mu\nu}u_\nu\\
 &={q^2\over6\pi\varepsilon_0c^5}
 \left(
 \dot a^\mu u^\nu u_\nu
 -u^\mu\dot a^\nu u_\nu
 \right).
\end{align*}

由 $u^2=c^2$ 一次微分得 $u\cdot a=0$,再微分一次得

\begin{equation*}
 a^2+u\cdot\dot a=0,
 \qquad
 u\cdot\dot a=-a^2.
\end{equation*}

因此有限自场四力化成

\begin{equation*}
 f_{\mathrm{self}}^\mu
 ={q^2\over6\pi\varepsilon_0c^3}
 \left(
 \dot a^\mu+{a^2\over c^2}u^\mu
 \right).
\end{equation*}

在号差 $+2$ 下,类时 $u^2=c^2$ 而类空 $a^2<0$。式中第二项必须写成加号,因为

\begin{align*}
 u_\mu f_{\mathrm{self}}^\mu
 &={q^2\over6\pi\varepsilon_0c^3}
 \left(u\cdot\dot a+{a^2\over c^2}u^2\right)\\
 &={q^2\over6\pi\varepsilon_0c^3}(-a^2+a^2)=0.
\end{align*}

这保证自力与四速度正交,不改变 $u^2=c^2$ 的归一化。

把正则场的分解与有限自力代回重整化方程,得到 Abraham--Lorentz--Dirac 方程

\begin{equation*}
 ma^\mu
 =f_{\mathrm{ext}}^\mu
 +{q^2\over6\pi\varepsilon_0c^3}
 \left(
 \dot a^\mu+{a^2\over c^2}u^\mu
 \right),
 \qquad
 f_{\mathrm{ext}}^\mu=qF_{\mathrm{in}}^{\mu\nu}u_\nu.
\end{equation*}

所以正则场在世界线上的有限极限、这里的有限自力和 ALD 方程不是三个独立结论。它们依次是同一个有限对象的场强形式、四力形式和运动方程形式。最后,世界管方法还可以给出同一自作用的第三种等价组织方式:把有限自力改写为束缚近场动量与远区辐射动量的守恒关系。

\subsubsection{束缚动量、Schott 项与远区辐射}

世界管方法还可以把有限自力改写成近场束缚动量和远区辐射动量。端帽内的电磁四动量 $P_{\mathrm{bound}}^\mu(\ell)$ 由奇异场在管内的积分 $(1/c)\int_{r\le\ell}T_{\mathrm{em}}^{\mu0}\d^3x$ 直接展开:Coulomb 部分给出 $q^2u^\mu/(8\pi\varepsilon_0c^2\ell)$(共动系中只有时间分量 $U_{\mathrm C}(\ell)/c$,提升后即与 $u^\mu$ 平行),世界线加速度的线性修正给出 $-q^2a^\mu/(6\pi\varepsilon_0c^3)$,故在小 $\ell$ 下

\begin{equation*}
 P_{\mathrm{bound}}^\mu(\ell)
 ={q^2\over8\pi\varepsilon_0c^2\ell}u^\mu
 -{q^2\over6\pi\varepsilon_0c^3}a^\mu
 +O(\ell).
\end{equation*}

下面两个性质将验证这一展开:第一项的导数复现管壁通量的发散部分,而总守恒式回到 ALD 方程。第一项是 Coulomb 近场携带的发散四动量,它的导数为

\begin{equation*}
 {\d\over\d\tau}
 \left({q^2\over8\pi\varepsilon_0c^2\ell}u^\mu\right)
 ={q^2\over8\pi\varepsilon_0c^2\ell}a^\mu,
\end{equation*}

正好复现管壁通量中的发散惯性。有限部分定义为 Schott 四动量

\begin{equation*}
 P_{\mathrm{Schott}}^\mu
 =-{q^2\over6\pi\varepsilon_0c^3}a^\mu.
\end{equation*}

它依赖瞬时加速度,是附着在粒子附近的可逆束缚场动量,而不是已经传播到无穷远的辐射动量。

远区 Lienard 辐射的四动量率必须与 $u^\mu$ 平行,因为在瞬时静止系中辐射的净三动量为零,而能量为正。协变形式为

\begin{equation*}
 {\d P_{\mathrm{rad}}^\mu\over\d\tau}
 =-{q^2\over6\pi\varepsilon_0c^5}a^2u^\mu.
\end{equation*}

由于 $a^2<0$,右侧时间分量为正。

把裸粒子动量、束缚场动量和已辐射动量放在同一个守恒式中:

\begin{equation*}
 {\d\over\d\tau}
 \left[m_0(\ell)u^\mu
 +P_{\mathrm{bound}}^\mu(\ell)\right]
 +{\d P_{\mathrm{rad}}^\mu\over\d\tau}
 =f_{\mathrm{ext}}^\mu.
\end{equation*}

逐项代入束缚动量和辐射四动量率的展开:

\begin{align*}
 &\left[m_0(\ell)
 +{q^2\over8\pi\varepsilon_0c^2\ell}\right]a^\mu
 -{q^2\over6\pi\varepsilon_0c^3}\dot a^\mu
 -{q^2\over6\pi\varepsilon_0c^5}a^2u^\mu
 +O(\ell)
 =f_{\mathrm{ext}}^\mu.
\end{align*}

使用质量重整化定义并把后两项移到右边,立即回到 ALD 方程。因此 Green 函数半差法、世界管壁通量法以及"裸粒子加束缚近场加远区辐射"的守恒法,是同一个自作用问题的三种等价组织方式。

具体地,由 Schott 四动量与远区辐射四动量率,

\begin{align*}
 -{\d P_{\mathrm{Schott}}^\mu\over\d\tau}
 -{\d P_{\mathrm{rad}}^\mu\over\d\tau}
 &={q^2\over6\pi\varepsilon_0c^3}\dot a^\mu
 +{q^2\over6\pi\varepsilon_0c^5}a^2u^\mu\\
 &={q^2\over6\pi\varepsilon_0c^3}
 \left(\dot a^\mu+{a^2\over c^2}u^\mu\right)
 =f_{\mathrm{self}}^\mu.
\end{align*}

第一行把自力看成两个四动量流的负和,第二行把同一个量化成局域世界线导数。

辐射四动量率的时间分量给出 Lienard 功率。四加速度不变量为

\begin{equation*}
 -a^2
 =\gamma^6
 \left[
 \vb a^2-{(\vb v\times\vb a)^2\over c^2}
 \right].
\end{equation*}

这一恒等式来自 $a^\mu$ 的三维分解:$a^0=\gamma^4\vb v\cdot\vb a/c$、$\vb a_{(4)}=\gamma^2\vb a+\gamma^4(\vb v\cdot\vb a)\vb v/c^2$,代入 $-a^2=(a^0)^2-\vb a_{(4)}^2$ 并合并同类项即得。故

\begin{equation*}
 P_{\mathrm{rad}}
 ={q^2\gamma^6\over6\pi\varepsilon_0c^3}
 \left[
 \vb a^2-{(\vb v\times\vb a)^2\over c^2}
 \right].
\end{equation*}

确实,由 $P^0=E/c$ 与 $u^0=\gamma c$,功率 $P=\d E/\d t=(1/\gamma)\d E/\d\tau=(c/\gamma)\d P_{\mathrm{rad}}^0/\d\tau=-q^2a^2/(6\pi\varepsilon_0c^3)$,代入上式一致。

在非相对论极限,定义辐射反作用时间

\begin{equation*}
 \tau_{\mathrm{rr}}
 ={q^2\over6\pi\varepsilon_0mc^3},
 \qquad
 m\tau_{\mathrm{rr}}={q^2\over6\pi\varepsilon_0c^3}.
\end{equation*}

则

\begin{equation*}
 \vb F_{\mathrm{AL}}=m\tau_{\mathrm{rr}}\dot{\vb a},
 \qquad
 P_{\mathrm{rad}}=m\tau_{\mathrm{rr}}\vb a^2.
\end{equation*}

机械功率由同一个自力给出

\begin{align*}
 \vb F_{\mathrm{AL}}\cdot\vb v
 &=m\tau_{\mathrm{rr}}\dot{\vb a}\cdot\vb v\\
 &=m\tau_{\mathrm{rr}}{\d\over\d t}(\vb a\cdot\vb v)
 -m\tau_{\mathrm{rr}}\vb a^2\\
 &=-{\d E_{\mathrm{Schott}}\over\d t}-P_{\mathrm{rad}},
 \qquad
 E_{\mathrm{Schott}}=-m\tau_{\mathrm{rr}}\vb a\cdot\vb v.
\end{align*}

所以瞬时机械功率并不总等于负的远区辐射功率;二者之差是束缚近场能量的时间导数。只有当散射过程初末 $\vb a\to0$ 时,Schott 边界项才消失,机械总能量损失才等于远区全频谱积分。

作为检验,考虑一维恒定固有加速度 $\alpha$ 的双曲运动。此时

\begin{equation*}
 a^2=-\alpha^2,
 \qquad
 \dot a^\mu={\alpha^2\over c^2}u^\mu
 =-{a^2\over c^2}u^\mu.
\end{equation*}

代入有限自力得

\begin{equation*}
 f_{\mathrm{self}}^\mu=0,
\end{equation*}

而 Lienard 功率仍给出非零辐射。由前面的链式关系,此时

\begin{equation*}
 {\d P_{\mathrm{Schott}}^\mu\over\d\tau}
 =-{\d P_{\mathrm{rad}}^\mu\over\d\tau},
\end{equation*}

即远区辐射四动量完全由束缚近场四动量的减少供给;这不是能量守恒的矛盾,而是瞬时场能分组的结果。

把 ALD 方程写成三维形式还可检验低速极限。四力与三力满足

\begin{equation*}
 f^\mu=\gamma\left({\vb F\cdot\vb v\over c},\vb F\right),
\end{equation*}

故三力是四力空间分量除以 $\gamma$。由 $u^\mu=\gamma(c,\vb v)$ 连续求导,

\begin{align*}
 a^0&=\gamma^4{\vb v\cdot\vb a\over c},\\
 \vb a_{(4)}&=\gamma^2\vb a
 +\gamma^4{\vb v\cdot\vb a\over c^2}\vb v,
\end{align*}

再对固有时求导得四 jerk 的各个分量,连同 $a^2u^\mu/c^2$ 项代入有限自力的空间分量 $\vb f_{\mathrm{self}}=(q^2/6\pi\varepsilon_0c^3)(\dot{\vb a}_{(4)}+a^2\gamma\vb v/c^2)$,除以 $\gamma$ 并合并同类项,得到相对论三维 ALD 力

\begin{align*}
 \vb F_{\mathrm{ALD}}
 ={q^2\over6\pi\varepsilon_0c^3}
 \bigg[{}
 &\gamma^2\dot{\vb a}
 +{3\gamma^4\over c^2}(\vb v\cdot\vb a)\vb a
 +{\gamma^4\over c^2}(\vb v\cdot\dot{\vb a})\vb v\\
 &+{3\gamma^6\over c^4}(\vb v\cdot\vb a)^2\vb v
 \bigg].
\end{align*}

当 $v/c\ll1$ 时,$\gamma=1+O(v^2/c^2)$,后三项均比第一项高两个 $1/c$ 阶,故

\begin{equation*}
 \vb F_{\mathrm{AL}}
 ={q^2\over6\pi\varepsilon_0c^3}\dot{\vb a}
 +O(c^{-5}).
\end{equation*}

因此非相对论 Abraham--Lorentz 力不是另一个独立模型,而是 Dirac 正则场自力在低速极限下的同一表达。
